\section{Conclusion}

This paper frames Context Sphere as a falsifiable architecture for topology-aware context allocation in software-engineering agents. In the current 100-case SWE-bench Verified ablation, Context Sphere achieved 41 local \texttt{PASS\_TO\_PASS} repair passes, compared with 36 for dense Vector retrieval and 39 for a naive Hybrid retriever, at roughly \$0.08 per issue. A follow-on projection smoke test shows that persona-conditioned routing can preserve 9/10 known Context Sphere successes while reducing input-token load by 71.5\% and estimated inference cost by 58.4\%.

The ablation also exposes the next research target. While Context Sphere provides a structured topological map of the repository, the 51-pass overlap ceiling shows that no single retrieval policy captures all solvable cases. The Context Projection Model is our first step toward that target: a trained discriminative routing layer designed to slice the Master Context Sphere into token-efficient, persona-specific \textit{Spheres of Influence} for the Product Manager, Worker, and Reviewer. The next stage is to scale this projection layer beyond the preliminary smoke cohort, pair it with stronger patch materialization, and evaluate the full stack under official SWE-bench \texttt{FAIL\_TO\_PASS} conditions.
